\documentclass[11pt]{letter}
\usepackage[margin=1in]{geometry}
\usepackage{hyperref}

% Signed by the corresponding author on behalf of both, which is the usual
% convention: the corresponding author is the one in correspondence with the
% journal. The co-author is named so the representation is explicit rather
% than left to be inferred from the body.
\signature{Sumanshu Sohal\\
\emph{on behalf of both authors}\\[6pt]
Department of Information Technology\\
University of the Cumberlands\\
Williamsburg, KY 40769, USA\\
sumanshu.95s@outlook.com}

\address{Sumanshu Sohal\\
Department of Information Technology\\
University of the Cumberlands\\
Williamsburg, KY 40769, USA}

\date{\today}

\begin{document}

\begin{letter}{Editors\\
Empirical Software Engineering\\
Springer}

\opening{Dear Editors,}

We submit our manuscript ``Cross-Corpus Evaluation of Domain-Aware Query
Reformulation for Regulatory Traceability'' for consideration in
\emph{Empirical Software Engineering}.

The paper's contribution is a way of reporting retrieval effects that shows
how much they can be trusted, applied to a case where the answer is
uncomfortable. Compliance mapping has few public collections, so methods are
usually measured on benchmarks their own authors build. We built one, measured
a $+0.121$ Top-1 gain from domain-aware query reformulation on it, and then
asked whether that measurement carried. Under a single ranking protocol
across four corpora it did not: three externally authored corpora give
$+0.028$, $0.000$ and $+0.032$, none distinguishable from zero.

Two properties of those estimates are the transferable part. Each turns on
fewer than ten requirements whose top-1 correctness changes at all, which no
confidence interval conveys, and each moves by up to $0.029$ when a
latent-semantic backend is refitted on a different population, a choice
unrelated to the intervention. We therefore report wins, losses and ties with
every paired effect, and a three-cell factorial separating evaluation protocol
from fitting population.

The work fits EMSE on three counts: it is an empirical study in requirements
engineering and traceability; it tests generalization under one common
protocol rather than comparing numbers produced under different ones; and it
ships infrastructure others can extend, namely four corpora, a frozen analysis
specification, the evaluation harness, and provenance tooling that checks
every number in the paper against a hashed record. One analysis, a gated
lexical--semantic hybrid, was preregistered before it was run
(doi:10.17605/OSF.IO/NZXRV); it established non-inferiority without
superiority while producing $+0.085$ on one corpus against $-0.088$ on
another. The corpora accompany the article as Online Resource~1 and the code
is public.

We are explicit about the limits. The three external corpora share a target
catalogue and NIST-authored crosswalks, so they are correlated evidence rather
than independent samples; with one corpus of our own the design cannot say
which of its properties matters; and our vocabulary-gap measure is fitted per
corpus, so it cannot rank corpora. We make no claim that author-built
benchmarks inflate results in general. An earlier unpublished two-corpus
analysis motivated this study; the additional corpora and common evaluation
protocol did not support its conditional vocabulary-gap interpretation.

The manuscript is original and is not under consideration elsewhere. Both
authors, Sumanshu Sohal and Darshankumar Prajapati, have read and approved the
submission, and we declare no competing interests.

\closing{Yours sincerely,}

\end{letter}
\end{document}
